\section{Hipóteses de planejamento}
A Tabela~\ref{tab:hipoteses} resume os parâmetros adotados. Alguns vêm de normas, como os da PNAB
\citep{pnab2017}, outros de valores de referência de custo, como os do Novo PAC Saúde
\citep{novopac2025}, e os demais são hipóteses do estudo. A análise de sensibilidade varia os raios, o crescimento populacional, o mínimo rural e a folga de custo, e os demais parâmetros ficam fixos. Os
parâmetros da PNAB são aplicados de modo uniforme, sem diferenciar a vulnerabilidade social ou as
condições locais de cada território.

\begin{table}[htbp]\centering
\caption{Parâmetros de planejamento e sua origem}\label{tab:hipoteses}
\begin{tabular}{p{.26\linewidth}p{.19\linewidth}p{.44\linewidth}}
\toprule Parâmetro & Valor principal & Origem e observação\\\midrule
Capacidade por eSF & 3.000 habitantes & Média recomendada desde a PNAB de 2011 e dentro da faixa de 2.000 a 3.500 da PNAB atual \citep{brasil2013teto,pnab2017}\\
Habitantes por ACS & 750 & Máximo da PNAB para áreas de maior dispersão ou vulnerabilidade, aplicado a todo o município\\
Raio de cobertura & 3 km & Hipótese do estudo, pois não há padrão oficial de distância. Os valores de 2 e 5 km são testados\\
Distância & Linha reta & Fórmula de Haversine, que dispensa uma base de vias e de tempos de viagem\\
Distância mínima entre candidata e UBS existente & 1 km & Hipótese do estudo, para que a nova UBS não divida a mesma área de uma existente\\
Mínimo para abrir UBS urbana & 2.000 habitantes & Limite inferior da faixa da PNAB, para que a nova equipe tenha população suficiente\\
Mínimo para abrir UBS rural & 500 habitantes & Hipótese do estudo, para permitir UBS em setores rurais pequenos. Os valores 0 e 1.000 são testados\\
Equipes por nova UBS & De 1 a 4 & Teto da PNAB e portes I a IV do Novo PAC Saúde\\
Folga operacional & 2\% & Hipótese do estudo, que permite reduzir a distância com pouco aumento de custo. O valor de 5\% é testado\\
Horizonte de custo & 5 anos & Hipótese do estudo, para somar construção e operação na mesma base, a preços constantes e sem desconto\\
Parcela da alta complexidade via secundária & 25, 50 e 75\% & Cenários, pois não há dados de encaminhamento por paciente\\
Capacidade da equipe móvel & 2.000 habitantes & Hipótese do estudo, igual ao limite inferior da faixa da PNAB por equipe\\
Alcance da equipe móvel & 15 km e 5 pontos & Hipótese do estudo, sem parâmetro oficial para equipes móveis\\
Peso da distância no índice & $\beta=1$ e $2$ & Sensibilidade, pois exige calibração empírica \citep{guagliardo2004}\\
\bottomrule\end{tabular}
\fonte{Elaboração própria, com base em \citet{pnab2017} e \citet{novopac2025}.}
\end{table}

\begin{table}[htbp]\centering
\caption{Origem dos parâmetros monetários e de pessoal}\label{tab:origem_custos}
\begin{tabular}{p{.26\linewidth}p{.23\linewidth}p{.40\linewidth}}
\toprule Parâmetro & Valor adotado & Origem e cálculo\\\midrule
Médico & R\$ 14.000/mês & Próximo da bolsa-formação bruta do Mais Médicos, de R\$ 14.121,63 em 2026 \citep{maismedicos2026}.\\
Enfermeiro & R\$ 5.000/mês & Hipótese do estudo, pouco acima do piso nacional de R\$ 4.750 \citep{brasil2022piso}.\\
Técnico de enfermagem & R\$ 3.325/mês & Piso nacional, 70\% do piso do enfermeiro, pelo valor integral \citep{brasil2022piso}.\\
Auxiliar de enfermagem & R\$ 2.375/mês & Piso nacional, 50\% do piso do enfermeiro, pelo valor integral \citep{brasil2022piso}.\\
ACS & R\$ 3.242/mês & Piso constitucional de dois salários mínimos de 2026, de R\$ 1.621 cada \citep{brasil2022,brasil2025sm}.\\
Provisões anuais & $(13+1/3)/12$ & Décimo terceiro e um terço de férias distribuídos nos 12 meses \citep{brasil1962decimo,brasil1988cf}.\\
Encargos sobre a base & 30\% & 20\% de contribuição patronal, 2\% de risco e 8\% de FGTS, aplicados a todos os vínculos \citep{brasil1991custeio,brasil1990fgts}.\\
Funcionamento sem pessoal & R\$ 12.000/mês por UBS existente e R\$ 15.000/mês por nova & Hipótese do estudo, sem dados municipais por unidade.\\
Construção & Conforme o porte & Valores do Novo PAC Saúde para o Sudeste (Tabela~\ref{tab:portes}), pagos uma única vez \citep{novopac2025}.\\
Coeficiente de deslocamento & R\$ 0,05/habitante-km & Hipótese do estudo, usada só para ordenar as soluções na terceira fase.\\
\bottomrule\end{tabular}
\fonte{Elaboração própria.}
\end{table}

\section{Índice de acessibilidade}
Além da cobertura, calcula-se um índice de acessibilidade que combina a oferta das equipes, a distância
e a população que disputa cada UBS. O índice segue a fórmula gravitacional de \citet{hansen1959}, com o
ajuste pela população concorrente de \citet{joseph1982}. Distâncias menores que 0,1 km são arredondadas
para esse valor, o que evita divisão por zero ($\widetilde d_{ij}$).
\begin{align}
V_j&=\sum_k\frac{P_k}{\widetilde d_{kj}^{\,\beta}}, &
A_i&=\sum_j\frac{S_j}{\widetilde d_{ij}^{\,\beta}V_j}, &S_j&=3000\,n_j.
\end{align}

\penalty0
O termo $V_j$ mede a população que disputa a UBS $j$, ponderada pela distância, e $A_i$ é a capacidade de
eSF potencialmente disponível por habitante do ponto $i$. O parâmetro $\beta$ controla o peso da
distância. O índice relativo divide $A_i$ pela média municipal, de modo que o valor 1 representa a média
da cidade. Ao contrário da cobertura, o índice considera todas as UBS, sem o corte de 3 km.

\section{Custos, orçamento e implementação}
O custo de construção das novas UBS segue os valores do Novo PAC Saúde para a região Sudeste, sem
terreno, equipamentos ou mobiliário (Tabela~\ref{tab:portes}) \citep{novopac2025}. O custo fixo de R\$ 15.000 por mês é o
mesmo para qualquer porte, o que pode subestimar a diferença de despesas entre unidades grandes e
pequenas. Como o número de equipes define o porte, as UBS com uma equipe são de porte I, e as de três ou
quatro equipes, de porte III ou IV.

\begin{table}[htbp]\centering
\caption{Custo de construção por porte de UBS}\label{tab:portes}
\begin{tabular}{lrr}\toprule Porte & Equipes & Construção (R\$)\\\midrule
I & 1 & 1.988.000\\II & 2 & 2.542.000\\III & 3 & 3.307.000\\IV & 4 & 4.306.000\\\bottomrule
\end{tabular}
\fonte{Manual Novo PAC Saúde, região Sudeste \citep{novopac2025}.}
\end{table}

Para a comparação com o orçamento, estima-se também o custo da rede atual nas mesmas categorias. O
pessoal clínico e de enfermagem, calculado pelas horas CNES e pelos custos unitários da
Tabela~\ref{tab:origem_custos}, custa cerca de R\$ 1,10 milhão por mês, e os 140 ACS, R\$ 0,66 milhão.
Com o custo fixo das 22 UBS, a rede atual custa R\$ 24,2 milhões por ano. A referência é a dotação
atualizada da Atenção Básica de 2026, isto é, o valor autorizado no orçamento municipal, de
R\$ 50,9 milhões, prevista na Lei Orçamentária Anual (LOA) e informada no Relatório Resumido da Execução
Orçamentária (RREO) \citep{itauna2025,tesouro2026}.

A dotação não é um saldo livre, pois também financia odontologia, medicamentos e administração, que não
entram no modelo. Por isso, a comparação indica o peso do plano no orçamento, e não a sua viabilidade
financeira. Como referência, o Sistema de Informações sobre Orçamentos Públicos em Saúde (SIOPS) registra
R\$ 32,5 milhões gastos com pessoal na Atenção Básica em 2025, valor que inclui categorias além das eSF e
dos ACS \citep{siops}.

O modelo da atenção primária foi implementado em Python, com a biblioteca PuLP e o solver HiGHS, usando
tolerância de otimalidade de 0,01\% e limite de 600 segundos por fase \citep{mitchell2011,huangfu2018}. O
modelo de encaminhamento foi escrito na linguagem MathProg e resolvido pelo GLPK \citep{glpk}. Testes
automatizados verificam a conservação da população e dos volumes, os limites de capacidade e a
consistência dos custos. A trilha de reprodução está no Apêndice~\ref{ap:reproducao}.
